\section{Mathematical setting}
\label{sec:setting}

Let $\Omega$ be a nonempty finite state space. A finite jump generator is specified by rates
\[
  k(x,y) \ge 0 \qquad (x\ne y),
\]
with escape rate
\[
  \lambda(x)=\sum_{y\ne x} k(x,y)
\]
and conservative generator matrix
\[
  Q(x,y)=
  \begin{cases}
    k(x,y), & x\ne y,\\
    -\lambda(x), & x=y.
  \end{cases}
\]
The formal structure \lean{FiniteJumpGenerator} packages the nonnegative rates, the induced escape rates and real generator, and the reference measures used below. Conservation is proved as the row-sum identity $\sum_y Q(x,y)=0$.

\subsection{Fixed-jump-count charts}

A path with exactly $n$ jumps is represented by a state sequence and a holding-time sequence,
\[
  \gamma=((x_0,\ldots,x_n),(\tau_0,\ldots,\tau_n)).
\]
The final holding interval $\tau_n$ records the no-jump remainder after the last jump. In Lean the type is
\[
  \lean{JumpPath Omega n}
  = (\lean{Fin (n+1) -> Omega})\times
    (\lean{Fin (n+1) -> NNReal}).
\]
All finite jump counts are assembled into the dependent sum
\[
  \lean{FullPath Omega}=\sum_{n:\mathbb{N}} \lean{JumpPath Omega n}.
\]
The development defines measurable terminal-state, jump-time, and right-continuous trajectory maps on this carrier.

For a fixed horizon $T\ge 0$, the holding times should satisfy $\sum_{i=0}^{n}\tau_i=T$. Restricting an ambient $(n+1)$-dimensional product Lebesgue measure to this equality slice would give the zero measure. The formalization therefore uses the free simplex
\[
  \Delta_n=\left\{u\in[0,1]^n:\sum_{i<n}u_i\le 1\right\},
\]
sets
\[
  \tau_i=T u_i \quad (i<n),
  \qquad
  \tau_n=T-\sum_{i<n}\tau_i,
\]
and proves
\[
  \operatorname{vol}(\Delta_n)=\frac{1}{n!}.
\]
The $T^n$ scaling of the free coordinates therefore gives the familiar simplex mass $T^n/n!$. Time reversal acts on the completed holding-time vector by an affine involution of the free chart, enabling a reusable change-of-variables theorem. The construction uses the free-coordinate product-volume convention; it does not claim the intrinsic Hausdorff normalization of the embedded equality slice.

\subsection{Sector density}

For a path started from a fixed state $x$, the $n$-jump density has the standard jump-process form
\begin{equation}
\label{eq:sector-density}
  \mathbf{1}_{\{x_0=x\}}
  \left(\prod_{i=0}^{n-1} k(x_i,x_{i+1})\right)
  \exp\left(-\sum_{i=0}^{n}\lambda(x_i)\tau_i\right).
\end{equation}
The density is integrated against counting measure on state sequences and the scaled free-simplex reference. The resulting measure is \lean{sectorLawFrom T x n}; its mass is \lean{sectorMassFrom T x n}. Equation~\eqref{eq:sector-density} is standard. The mechanized contribution is the construction, normalization, and subsequent identification of the pushforwards of the countable sector sum.

